% ===================================================================
%  TABELL Nr. 16 — De grousse Buttek. --- De Basar. D'Spillsaachen.
%  Traduction 2026 d'apres texto/io, controlee sur texto/fr, texto/de
%  et texto/nl. Consigne : voir texto/lb/10-tabell-01.tex.
%
%  DERNIER TABLEAU DE LA COLONNE, et il met sur UNE SEULE PAGE les
%  deux moitiés que les quinze precedents avaient rencontrees separement.
%
%  A l'alinea 3, la menagerie en carton. Tout y est germanique, et
%  compose sur place :
%    « Flantermaus »  la chauve-souris, la souris-qui-voltige ;
%    « Nilpäerd »     l'hippopotame, le cheval du Nil ;
%    « Walfësch »     la baleine ;  « Haifësch » le requin ;
%    « Séihond »      le phoque, le chien de mer ;
%    « Schildkröt »   la tortue ;   « Wandhond » le levrier ;
%    « Rendéier », « Nashorn », « Wollef », « Fuuss », « Af ».
%  Le francais, pour les memes betes, n'a que des mots simples et
%  opaques : chauve-souris mise a part, ni hippopotame, ni baleine,
%  ni requin, ni phoque, ni tortue ne se laissent demonter.
%
%  Aux alineas 9 a 17, le magasin lui-meme. Tout y est francais :
%  « Keess », « Client », « Comptabilitéit », « Étalage »,
%  « Vitrinn », « Rayon », « Rayonschef », « Encaisseur »,
%  « Corseten », « Caleçonen », « Parfumerie », « Portmonni »,
%  « Portefeuille », « Éventail », « Expeditioun », « Inspekter »,
%  « Arrestatioun », « Comptoir », « Fauteuil », « Parterre »,
%  « Chef d'orchestre », « Rideau », « Coffre-fort ».
%
%  LA LIGNE NE PASSE DONC NI ENTRE LE HAUT ET LE BAS, NI ENTRE
%  L'ANCIEN ET LE NOUVEAU : ELLE PASSE ENTRE CE QUI EXISTE SANS
%  GUICHET ET CE QUI N'EXISTE QUE PAR UN GUICHET. Une bete est la
%  qu'on la nomme ou non ; un rayon, une caisse, une expedition, une
%  arrestation ne sont rien d'autre que le geste d'une institution, et
%  l'institution est arrivee avec son nom francais dans la bouche.
%
%  C'EST LE CRITERE QUE LE TABLEAU 14 AVAIT TROUVE SOUS SA FORME
%  ETROITE — le registre des corporations predit la langue d'un
%  metier — et le 16 le donne sous sa forme large. Il vaut aussi pour
%  le tableau 2 (le corps allemand, le sport francais), pour le 6
%  (on mange en allemand, on sert en francais), pour le 11 (la
%  « Spuerkeess » et les « Pompjeeën »).
%
%  MAIS IL PREDIT, IL NE DECIDE PAS, et le tableau 15 en donnait deja
%  le contre-exemple : le « Schwäinsmetzler » tient boutique, paie
%  patente et a eu sa corporation, et son nom est germanique de bout
%  en bout. L'armee des alineas 4 et 5 penche du cote francais
%  (« Arméi », « Assaut », « Attack », « Escadron », « Retraite »,
%  « Victoire », « Bataillon », « Carré ») mais garde ses « Zaldoten »,
%  son « Kampf », ses « Zelter », son « Waffestëllstand » et son
%  « Friddensvertrag ». LE CRITERE DIT OU CHERCHER ; IL NE DIT PAS CE
%  QU'ON TROUVERA.
%
%  L'ECART DECLARE DU BLOC t16-15-1, LE TROISIEME ET DERNIER DES
%  TROIS. Le fac-simile ido y ecrit « mi-sferi d) », sans la
%  parenthese ouvrante. Comme au tableau 8, la coquille est du
%  compositeur de l'ido et non du texte luxembourgeois : la parenthese
%  est retablie, comme dans les quatorze colonnes precedentes, et
%  renvoji.py passe l'ecart sans rien dire. UNE TRADUCTION N'EST PAS
%  UNE TRANSCRIPTION.
% ===================================================================

%%K t16-apar-1 apar
\VUsaut{4.0mm}
\VUtitre{15.0pt}{60}{TABELL Nr. 16}
\VUsaut{2.0mm}
\VUpk{13.2pt}{40}{De grousse Buttek. --- De Basar.}
\VUpk{13.2pt}{40}{D'Spillsaachen.}
\VUsaut{3.4mm}

%%K t16-01-1 p
1. --- Dat neit Joer kënnt no. Déi grouss Butteker hunn hir Étalage vu
\VUgras{Spillsaachen} \textsuperscript{(1)} a vun Neijoerschcadeauen preparéiert.

%%K t16-01-2 p
Ech hunn de Grousspapp gefrot, hie sollt mech an de Basar féieren, fir déi schéin
Ausstellung ze gesinn, an hien huet jo gesot.

%%K t16-02-1 p
2. --- Als éischt si mir virun schéine Waffeschëlder stoe bliwwen, op deenen e
\VUgras{Gewier}, e \VUgras{Kepi}, e \VUgras{Sabel}, en \VUgras{Ceinturon} an e Paar
\VUgras{Epauletten} sinn, oder e \VUgras{Helm}, e \VUgras{Kürass}
\textsuperscript{(3)} an eng \VUgras{Pistoul} \textsuperscript{(4)}. Dat alles huet
mech vill méi interesséiert wéi déi \VUgras{Zauberlatären} \textsuperscript{(2)}.

%%K t16-tit-1 sub
\VUsaut{3.0mm}
\VUpk{10.2pt}{40}{Schéin Ubléck.}
\VUsaut{2.6mm}

%%K t16-03-1 p
3. --- An deemselwechte Rayon stoung e \VUgras{Schëldwaach} \textsuperscript{(5)} a
sengem \VUgras{Waachthaischen}; et huet ausgesinn, wéi wann hie virun enger ganzer
Menagerie aus Karton Wuecht géif halen. Wéi vill Déiere waren do: e \VUgras{Papagei}
\textsuperscript{(6)} op senger Sëtzstaang, en \VUgras{Nashorn} \textsuperscript{(7)}
mat engem Horn op der Nues, e \VUgras{Rendéier} \textsuperscript{(8)}, e
\VUgras{Strauss} \textsuperscript{(9)}, en \VUgras{Af} \textsuperscript{(10)}, deen eng
Noss knuppert, e \VUgras{Fuuss} \textsuperscript{(11)}, deen eng Hong dovundréit, e
\VUgras{Krokodil} \textsuperscript{(12)}, dee seng immens Kiefer opräisst, e
\VUgras{Kamel} \textsuperscript{(13)}, en elegante \VUgras{Wandhond}
\textsuperscript{(14)}, e \VUgras{Panther} \textsuperscript{(15)}, duerno zwee
Déieren, déi ech net kannt hunn an déi, wéi ee seet, e \VUgras{Wisel}
\textsuperscript{(16)} an eng \VUgras{Hyän} \textsuperscript{(17)} sinn. Och do waren e
\VUgras{Leopard} \textsuperscript{(18)} an e \VUgras{Wollef} \textsuperscript{(21)};
ganz nobäi waren zaart \VUgras{Rätercher} \textsuperscript{(19)} a winzeg
\VUgras{Mailercher} \textsuperscript{(20)} aus Gummi. Als Lescht hunn en
\VUgras{Nilpäerd} \textsuperscript{(22)} an e \VUgras{Dromedar} \textsuperscript{(23)}
mat sengem Bockel d'Sammlung ofgeschloss. Ech wier do nach vill Stonne bliwwen, wann
net niewendrun e prächtegt Lager voller \VUgras{Bläizaldoten} \textsuperscript{(29)}
gewiescht wier, dat meng Opmierksamkeet ugezunn huet.

%%K t16-tit-2 sub
\VUsaut{4.0mm}
\VUpk{10.2pt}{40}{Zaldoten a Krich.}
\VUsaut{2.8mm}

%%K t16-04-1 p
4. --- Mäi Grousspapp, deen \VUgras{Invalid} \textsuperscript{(24)} ass an deen
d'Arméi huet misse verloossen wéinst enger \VUgras{Wonn,} déi him
d'\VUgras{Militärmedail} agebruecht huet, erzielt mir ganz gär vun de
\VUgras{Campagnen,} un deenen hien deelgeholl huet. Hie war frou, mir déi
verschidde Beweegunge vun der klenger Miniaturarméi z'erklären. Eng Grupp vu
richtege Zaldoten, déi de Basar besicht hunn --- e \VUgras{Genie-Zaldot}
\textsuperscript{(25)}, en \VUgras{Alpejeeër} \textsuperscript{(26)} mat engem
\VUgras{Béret} um Kapp, e \VUgras{Sergent-Major} \textsuperscript{(27)} vun der
Infanterie an e \VUgras{Maréchal des logis} \textsuperscript{(28)} vun der Artillerie
mat engem gëllene \VUgras{Galon} um Aarm ---, si bei eis stoe bliwwen, fir
d'Erklärungen unzelauschteren.

%%K t16-05-1 p
5. --- Mäi Grousspapp, ganz stolz, datt hien esou e brillanten Publikum hat, huet e
ganze Schluechtplang improviséiert.

%%K t16-05-2 p
Déi festungsstark Stad mat hire \VUgras{Wäll} an hire \VUgras{Griewen} gouf vun der
\VUgras{feindlecher Arméi} belagert, déi no engem laange \VUgras{Bombardement} prett
war fir den \VUgras{Assaut.} Mä déi \VUgras{Belagert,} vum Honger gedriwwen, hunn en
\VUgras{Ausfall} op de \VUgras{Camp} vun de \VUgras{Belagerer} probéiert. Nodeem si
d'\VUgras{Attack} an engem verbësse \VUgras{Kampf} ronderëm d'\VUgras{Zelter}
zeréckgeschloen haten, hunn déi \VUgras{gewonnen} Belagerer hir \VUgras{Escadronen}
hannert den \VUgras{geschloenen} Ugräifer hier geheit. Déi hunn d'\VUgras{Retraite}
ugetrueden an d'\VUgras{Schluechtfeld} mat \VUgras{Doudegen} a \VUgras{Blesséierte}
bedeckt. \VUgras{Brancardieren} hunn déi Lescht an d'\VUgras{Ambulanz} gedroen, fir
se vum \VUgras{Chirurg} versuerge ze loossen.

%%K t16-05-3 p
Trotz dem heldenhafte Widderstand vun e puer \VUgras{Bataillonen,} déi e
\VUgras{Carré} gebilt haten, war d'\VUgras{Néierlag} vollstänneg, de Rescht vun der
Arméi huet d'\VUgras{Flucht} ergraff an huet vill \VUgras{Gefaangen} hannerlooss. De
Feind ass gaangen, seng \VUgras{Victoire} ze vollenden, andeem hien d'Stad besat huet.
Vläicht wier dat d'Enn vun der Campagne, an no engem \VUgras{Waffestëllstand} kéint
een e \VUgras{Friddensvertrag} ënnerschreiwen.

%%K t16-tit-3 sub
\VUsaut{4.4mm}
\VUpk{10.2pt}{40}{Neijoerschcadeauen.}
\VUsaut{2.8mm}

%%K t16-06-1 p
6. --- Déi jonk Zaldoten, déi beim Nolauschtere vun der faarweger Erzielung vum
Veteran gelaacht hunn, hunn hien nach e puer aner Erklärunge gefrot. Ech fir mäin
Deel sinn ronderëm de Comptoir gaangen, fir eng schéin \VUgras{Aarmbrëscht}
\textsuperscript{(32)} an e \VUgras{Bou} \textsuperscript{(33)} mat sengem
\VUgras{Päil} \textsuperscript{(31)} ze kucken. Do hunn ech eng aner Sammlung vun
Déiere fonnt: zwee \VUgras{Sangvullen} \textsuperscript{(34)} --- eng automatesch
\VUgras{Nuechtigall} an eng \VUgras{Grasmécks,} déi aus voller Kiel sangen --- an eng
Rei komesch Déieren: eng \VUgras{Flantermaus} \textsuperscript{(35)}, eng \VUgras{Boa}
\textsuperscript{(36)}, eng \VUgras{Schildkröt} \textsuperscript{(37)}, e
\VUgras{Séihond} \textsuperscript{(38)}, e \VUgras{Walfësch} \textsuperscript{(39)} an
en opgeblosene \VUgras{Haifësch} \textsuperscript{(40)}.

%%K t16-07-1 p
7. --- Ech si net laang stoe bliwwen, fir se ze bekucken, well ech dat gesinn hunn,
wourop ech gedreemt hunn: e prächtegt \VUgras{Marionettentheater} \textsuperscript{(41)}
mat wonnerschéinen \VUgras{Dekoren} an engem entzéckende roude \VUgras{Rideau.} Op der
\VUgras{Zeen} hunn zwee Acteuren den Androck gemaach, si géifen eng \VUgras{Roll} an
engem ganz interessante \VUgras{Stéck} spillen, well déi kleng \VUgras{Zuschauer} aus
Karton, an de Logen installéiert oder op de \VUgras{Fauteuilen} an am
\VUgras{Parterre} hannert dem \VUgras{Chef d'orchestre} sëtzend, liewhaft applaudéiert
hunn.

%%K t16-07-2 p
Leider huet dat Theater ganz deier gekascht.

%%K t16-08-1 p
8. --- Iwwregens war et schwéier, d'Spillsaachen auszewielen, well an de Vitrinne
Spillsaache vun alle Sorten opgestapelt waren: \VUgras{Arlequin} \textsuperscript{(42)},
\VUgras{Polichinelle} \textsuperscript{(43)}, \VUgras{Pierrot} \textsuperscript{(44)},
\VUgras{Eilen} \textsuperscript{(45)}, déi mat de Flilleke schloen, päifend
\VUgras{Märelen} \textsuperscript{(46)}, blaffend \VUgras{Pudelen}, e
\VUgras{Phonograph} \textsuperscript{(47)} a \VUgras{Gedëllegkeetsspiller.} Ee vun
deene Spillsaachen huet mech ganz amüséiert: et huet eng \VUgras{Séischluecht}
\textsuperscript{(48)} duergestallt, an där e \VUgras{Schëff} vu \VUgras{Piraten}
nieft engem Land ugegraff gouf, dat mat \VUgras{Kokosnossbeem} bepflanzt war.

%%K t16-noto-1 noto
\VUnotes{143.2mm}{%
\textit{(*) Kuck d'Tabell Nr.~6.}%
}

%%K t16-09-1 p
9. --- Ech konnt all déi Wonner ganz einfach bekucken, well nëmme wéineg Leit am
Buttek waren, kaum e puer Flâneuren, e \VUgras{Dragoun} \textsuperscript{(49)} an en
\VUgras{Encaisseur} \textsuperscript{(50)}, deen eng \VUgras{Tratte} bei der
\VUgras{Haaptkeess} \textsuperscript{(51)} presentéiert huet.

%%K t16-10-1 p
10. --- Ech hunn de Grousspapp gefrot, wat déi Bicher solle, déi op de Briedere
hannert der \VUgras{Keessiererin} louchen; hien huet mir geäntwert, datt et
\VUgras{Comptabilitéitsbicher} sinn.

%%K t16-tit-4 sub
\VUsaut{3.6mm}
\VUpk{10.2pt}{40}{Gafen ronderëm de Buttek.}
\VUsaut{2.8mm}

%%K t16-11-1 p
11. --- Wéi mir de Rayon vun de Spillsaache verlooss hunn, si mir an d'Suedlerei
gaangen, fir e Bléck op d'\VUgras{Suedelen} \textsuperscript{(53)},
d'\VUgras{Steigbiggelen} \textsuperscript{(54)}, d'\VUgras{Zeem} \textsuperscript{(55)},
d'\VUgras{Gebësser} \textsuperscript{(56)} an d'\VUgras{Cravachen} ze werfen. Wärend de
\VUgras{Rayonschef} \textsuperscript{(57)} mengem Grousspapp e puer Informatioune ginn
huet, sinn ech den Étalage vum \VUgras{Porzeläin} an de \VUgras{Kristall}
\textsuperscript{(58)} kucke gaangen, wou et schéin \VUgras{Zalotschosselen} a fein
\VUgras{Téikannen} \textsuperscript{(59)} gëtt.

%%K t16-12-1 p
12. --- Fir op den éischte Stack ze goen, si mir hannert der Vitrinn vun de
\VUgras{Corseten} \textsuperscript{(60)} laanschtgaangen, nieft der Bonneterie, wou
ganz schéi \VUgras{Caleçonen} \textsuperscript{(63)} ausgestallt waren. Niewendrun huet
eng \VUgras{Verkeeferin} \textsuperscript{(61)} eng \VUgras{Keeferin}
\textsuperscript{(62)} schéi seiden \VUgras{Stoffer} \textsuperscript{(64)} auswiele
gelooss.

%%K t16-12-2 p
Beim Erklamme vun der Trap hunn ech gemierkt, datt de Buttek mat japanesche
\VUgras{Latären} \textsuperscript{(65)}, \VUgras{Parasolen} \textsuperscript{(66)} an
immens grousse \VUgras{Palmen} \textsuperscript{(70)} dekoréiert ass.

%%K t16-13-1 p
13. --- Um éischte Stack war net vill ze gesinn; nëmme Miwwelen (*),
\VUgras{Coffres-forts} \textsuperscript{(67)}, \VUgras{Kommoden} \textsuperscript{(68)},
\VUgras{Kannerween} \textsuperscript{(69)}. Mä de Gesamtubléck vum Buttek war ganz
\VUgras{amüsant.}

%%K t16-14-1 p
14. --- Bei eise Féiss hunn ech d'Museksinstrumenter gesinn: \VUgras{Harfen}
\textsuperscript{(71)}, \VUgras{Gittaren} \textsuperscript{(72)},
\VUgras{Mandolinnen} \textsuperscript{(73)}, \VUgras{Piston-Cornetten}
\textsuperscript{(74)}, \VUgras{Klarinetten} \textsuperscript{(75)}, Geien mat hirem
\VUgras{Streechbou} \textsuperscript{(76)} a souguer eng \VUgras{grouss Tromm}
\textsuperscript{(77)}. E bësse méi wäit ewech, um Rayon vum Pabeier, huet e
\VUgras{Student} \textsuperscript{(78)} mat dem \VUgras{Verkeefer}
\textsuperscript{(79)} ëm eng Schoulmapp gefeilscht. Nieft hinnen hunn ech e schéint
\VUgras{ABC-Buch} \textsuperscript{(80)} erkannt.

%%K t16-14-2 p
An der Mëtt vum Buttek stoung en héije \VUgras{Chrëschtbam} \textsuperscript{(81)},
ganz mat Käerzen, Klengegkeeten a Spillsaache vun alle Sorte behaangen:
\VUgras{Holzpäerd} \textsuperscript{(82)}, \VUgras{Poppen} \textsuperscript{(83)}, asw.

%%K t16-14-3 p
D'\VUgras{Parfumerie} \textsuperscript{(84)} mat hire \VUgras{Zännwaasser} an hire
verschiddene \VUgras{Parfumen} huet e ganz gudde Geroch verbreet, deen bis bei eis
eropgestige ass.

%%K t16-tit-5 sub
\VUsaut{3.4mm}
\VUpk{10.2pt}{40}{Mir kafen eppes.}
\VUsaut{2.8mm}

%%K t16-15-1 p
15. --- Well mäi Grousspapp d'Zäit ugefaangen huet ze laang ze fannen, si mir iwwer
déi grouss Trap erofgaangen. Hien ass e Moment virun enger \VUgras{Astronomietafel}
\textsuperscript{(85)} stoe bliwwen, déi, wéi hie mir gesot huet, \VUgras{Kometen}
\textsuperscript{(a)}, \VUgras{Mound- a Sonneklipsen} \textsuperscript{(b)}, eng
Wandrous mat de \VUgras{véier Himmelsrichtungen} \textsuperscript{(c)}
\VUgras{(Norden, Süden, Osten a Westen)}, eng Kaart vun deenen zwou
\VUgras{Hallefkugelen} \textsuperscript{(d)}, d'\VUgras{Planéiten}
\textsuperscript{(e)} an d'\VUgras{Sonnesystem} \textsuperscript{(f)} duergestallt
huet. Ech hunn näischt vun alldeem verstanen, a mech huet d'galonéiert Livrée vun
engem \VUgras{Liwwerjong} \textsuperscript{(86)}, dee laanschtgaangen ass, vill méi
interesséiert, an och déi schéin \VUgras{Éventailen} \textsuperscript{(87)}, déi nieft
mir waren, bei de \VUgras{Portmonnien} \textsuperscript{(88)} an de
\VUgras{Portefeuillen} \textsuperscript{(89)}.

%%K t16-16-1 p
16. --- Ech hunn och um Comptoir d'\VUgras{Fiederen} \textsuperscript{(90)} an
d'\VUgras{Rimenschnällen} \textsuperscript{(91)} bewonnert, déi schéi
\VUgras{kënschtlech Blummen} \textsuperscript{(94)}, déi graziéis a Kierf geluecht
waren. D'\VUgras{Vijoletten}, de \VUgras{Goldlack}, d'\VUgras{Hyazinthen}
\textsuperscript{(95)}, d'\VUgras{Geranien} \textsuperscript{(96)},
d'\VUgras{Lilien} \textsuperscript{(97)} an d'\VUgras{Kapuzinerblummen}
\textsuperscript{(98)} ware ganz gutt nogemaach.

%%K t16-tit-6 sub
\VUsaut{4.4mm}
\VUpk{10.2pt}{40}{De Cadeau vu mengem Grousspapp.}
\VUsaut{2.8mm}

%%K t16-17-1 p
17. --- Ech hunn de Grousspapp gebiet, mir eng vun den \VUgras{Zännbiischten}
\textsuperscript{(92)} ze kafen, déi an enger Boîte nieft den \VUgras{Hoernolen}
\textsuperscript{(93)} louchen. Hien huet jo gesot, an ech si selwer mäin Akaf bei der
Keess bezuele gaangen, bei där eng wichteg \VUgras{Expeditioun} \textsuperscript{(100)}
fir e \VUgras{Client} \textsuperscript{(99)} preparéiert gouf.

%%K t16-18-1 p
18. --- Op eemol ass eng liicht Onrou ënner de Leit entstanen, déi bei de
kënschtlereschen \VUgras{Bronzen} \textsuperscript{(101)} stoe bliwwe waren. En
\VUgras{Déif} \textsuperscript{(102)} war grad vun engem \VUgras{Inspekter}
\textsuperscript{(103)} op frëscher Dot erwëscht ginn, wéi hie probéiert huet, een
auszerauben. Ee gouf geschéckt, fir e Polizist ze huelen, deen d'\VUgras{Arrestatioun}
maache sollt, a grad wéi mir erausgaangen sinn, gouf de Verbriecher an d'Prisong
gefouert.

\VUsaut{6.0mm}
\VUfilet{22mm}
